\section{Necessity of section lifting for inseparability}
\label{app:lifting-counterexample}

The lifting hypothesis in Theorem~\ref{thm:local-lifting} cannot be
omitted, even for a faithful range functor that preserves restriction
products and reflects zero arrows.

Fix an acceptable numbering $(\varphi_e)_{e\in\N}$ of the ordinary
partial recursive functions $\N\rightharpoonup\N$. For a partial
numerical function $f$, write $f(x)\mathord\downarrow$ when $f$ is
defined at $x$. Put
\[
 K=\{e:\varphi_e(e)\mathord\downarrow\}.
\]
The set $K$ is not recursive.
Let $(W_a^K)_{a\in\N}$ enumerate the domains of an acceptable
numbering of partial functions recursive relative to $K$.
Fix a recursive bijective pairing function
$\langle-,-\rangle:\N^2\to\N$, with projections $\pi_0,\pi_1$.

The objects of $\mathcal C_K$ are the numerical carriers generated
from $\N$ by binary Cartesian products and tagged disjoint unions,
with fixed recursive codings. An arrow $f:X\rightharpoonup Y$ is a partial function whose
domain $D_f$ is c.e.\ relative to $K$ and which has an ordinary partial
recursive extension $f_0:X\rightharpoonup Y$. Thus $f_0$ is defined
on every element of $D_f$ and agrees with $f$ there. Equality is equality
of partial-function graphs. These arrows need not themselves be ordinary
partial recursive functions. The total arrows of $\mathcal C_K$ are exactly
the ordinary total recursive functions between the indicated types.

Let $F:\mathcal C_K\to\Pfn$ send each numerical carrier to the same
set and each arrow to the same partial function.

On $\N$, the arrows are precisely the functions
\[
 f_{e,a}(x)\simeq
 \begin{cases}
  \varphi_e(x),&x\in W_a^K,\\
  \text{undefined},&x\notin W_a^K.
 \end{cases}
\]
Here $\simeq$ includes simultaneous undefinedness: the actual domain
is $W_a^K\cap\operatorname{dom}(\varphi_e)$. Every arrow has this
form by taking an index of its extension and an oracle index of its
specified domain.

\begin{proposition}\label{prop:missing-lift}
The category $\mathcal C_K$ is a prodominical recursion category with
binary restriction coproducts and ranges. The functor $F$ satisfies
every hypothesis of
Theorem~\ref{thm:local-lifting} except the lifting hypothesis.
For the Turing morphism and coproduct isomorphism constructed below,
the specified pair is separated by complementary restriction idempotents in
$\mathcal C_K$. Moreover, there are parallel arrows in $\mathcal C_K$
with disjoint restriction idempotents and no common upper bound in
the restriction order, although finite joins of restriction idempotents
exist.
\end{proposition}

\begin{proof}
Let $f:X\rightharpoonup Y$ and $g:Y\rightharpoonup Z$ have
ordinary partial recursive extensions $f_0,g_0$. Their composite has
extension $g_0f_0$ and domain
\[
 D_{gf}=D_f\cap f_0^{-1}(D_g).
\]
This domain is c.e.\ relative to $K$: enumerate $D_f$, compute the
corresponding $f_0$-values, and test for their appearance in the
enumeration of $D_g$, using dovetailing. Identities and associativity
are inherited from partial functions. Restriction is the partial
identity on $D_f$, which has the ordinary identity as an extension.
Empty partial functions give restriction zeros.

Cartesian products and tagged copairings preserve the arrow class:
their domains are respectively products and tagged unions of the given
domains, and their extensions are the corresponding ordinary partial
recursive programs. Consequently the set-theoretic restriction
products and coproducts have their usual universal properties in
$\mathcal C_K$. Product with a zero arrow is zero, and the chosen
recursive coding isomorphisms make the category isotypical.

The set $f(D_f)$ is c.e.\ relative to $K$: enumerate $D_f$ and run
$f_0$ on its elements. The partial identity on this image belongs to
$\mathcal C_K$ and is its range. Leastness and the agreement law
hold for partial functions: equality after $f$ gives equality on
every value in its image, including equality of definedness.
Ranges commute with restriction products because images of products
of partial functions are products of images. Restriction idempotents
are the partial identities on the subsets c.e.\ relative to $K$;
their order corresponds to inclusion, and their finite joins are
partial identities on unions.

Define $\tau_K:\N\times\N\rightharpoonup\N$ by
\[
 \tau_K(\langle e,a\rangle,x)\simeq f_{e,a}(x).
\]
Its domain is c.e.\ relative to $K$, uniformly in all three
variables, and it has the ordinary partial recursive extension
$(\langle e,a\rangle,x)\mapsto\varphi_e(x)$.
To prove its Turing property, let
$v:\N\times\N\rightharpoonup\N$ be an arrow, with ordinary
partial recursive extension $v_0:\N\times\N\rightharpoonup\N$ and
domain $D_v$ c.e.\ relative to $K$. Specializing the ordinary program
for $v_0$ gives an ordinary total recursive function $e:\N\to\N$ with
\[
 \varphi_{e(z)}(x)\simeq v_0(z,x).
\]
Specializing an oracle program for $D_v$ gives an ordinary total
recursive function $a:\N\to\N$ with
\[
 W_{a(z)}^K=\{x:(z,x)\in D_v\}.
\]
The latter specialization is ordinary recursive because it only
inserts the finite numeral for $z$ into program code; the oracle is
used when the resulting program runs. Thus the function
$r:\N\to\N$ given by $r(z)=\langle e(z),a(z)\rangle$ is a total
morphism of $\mathcal C_K$, and
$v=\tau_K(r\times1_{\N})$, with equality of domains and values.
The fixed recursive codings give the same argument for every parameter
object of $\mathcal C_K$. Hence $\tau_K$ is a Turing morphism.

The functor $F$ is faithful and preserves restrictions, zeros,
ranges, binary restriction products, and coproducts. It also reflects
zero arrows. Its target $\Pfn$ is split and locally connected:
restriction idempotents split through their subsets of definition,
and partial identities on singletons detect partial functions.
Partial injections have partial inverses, giving weak choice, and
inverse images preserve unions. Also $F(1_{\N})\ne0$.

Use the total recursive isomorphism
\[
 b:\N\longrightarrow\N+\N,\qquad
 b(2n)=i_0(n),\quad b(2n+1)=i_1(n).
\]
Put $d=\tau_K\Delta_{\N}$ and
$a_0=(bd)^*(1_{\N}+0_{\N})$,
$a_1=(bd)^*(0_{\N}+1_{\N})$.
For $q\in\N$, write
$p(q)\simeq\varphi_{\pi_0(q)}(q)$. These restriction idempotents
are the partial identities on the sets
\[
 \begin{aligned}
 D_{a_0}&=\{q:q\in W_{\pi_1(q)}^K,
                    \ p(q)\mathord\downarrow\text{ with an even value}\},\\
 D_{a_1}&=\{q:q\in W_{\pi_1(q)}^K,
                    \ p(q)\mathord\downarrow\text{ with an odd value}\}.
 \end{aligned}
\]
Let
\[
 A=\{q:p(q)\mathord\downarrow\text{ with an even value}\}.
\]
This is an ordinary c.e.\ set and is decidable relative to $K$.
For the latter assertion, uniformly construct a program that ignores
its input, simulates $p(q)$, and halts precisely when the output is
even; membership of its index in $K$ decides membership of $q$ in $A$.
Thus $A$ and $A^c$ are c.e.\ relative to $K$, and their partial
identities are complementary restriction idempotents in $\mathcal C_K$.
They separate the displayed pair because
$D_{a_0}\subseteq A$ and $D_{a_1}\subseteq A^c$.

The set $A$ is not recursive. Fix any index $a_*$. Choose an ordinary
total recursive function $h:\N\to\N$ such that $\varphi_{h(n)}$
ignores its input, simulates $\varphi_n(n)$, and returns $0$ if that
simulation halts. Then
\[
 n\in K\quad\Longleftrightarrow\quad
 \langle h(n),a_*\rangle\in A.
\]
For $B\subseteq\N$, write $1_B:\N\rightharpoonup\N$ for the
partial identity on $B$. The source arrow
$p_A=\nabla_{\N}(1_{A^c}+1_A):\N+\N\rightharpoonup\N$
is sent to the partial injection whose unique section
$s_A:\N\to\N+\N$ in $\Pfn$ is
\[
 s_A(n)=
 \begin{cases}
  i_1(n),&n\in A,\\
  i_0(n),&n\notin A.
 \end{cases}
\]
This section is total and not recursive, since its tag decides $A$.
If it had a preimage under $F$, that source arrow would
be total and hence ordinary total recursive, a contradiction.
This is the missing lift required by
Theorem~\ref{thm:local-lifting} for the complemented restriction idempotent $1_A$.

Finally, the arrows $i_1 1_A$ and $i_0 1_{A^c}$ belong to
$\mathcal C_K$: they have the total recursive extensions $i_1$
and $i_0$, respectively. Their domains are disjoint. Any common
upper bound would agree with the first arrow on $A$ and the second
on $A^c$. It would therefore be the total function $s_A$, which
is not a source arrow. Thus no such upper bound exists.
\end{proof}

